\magnification=\magstephalf
\input notesmac.tex
\advance\vsize by 1cm
\nopagenumbers

\centerline{\titlefont CZ3106 Symbolic Computing, Lab 3}
\bigskip
\line{Dr.~Wang Jian-Sheng \hfill For the weeks 10-22 Feb, due Saturday 22 Feb 03}
%----------------------------------------------------------------------- 
\problem Finish the code for Gaussian elimination method for solving
linear equations.  You can copy my code for the elimination part; 
you need to supply your code on the {\tt backsub[..]} for backward 
substitution.  Test your program with at least one numerical input, and
one symbolic input. 

\problem Define a least-square fit function similar to the built-in
function {\tt Fit[]}.  Our fitting function takes the form\par
\ll{\t leastsquarefit[\char`\{\char`\{$x_1$,$y_1$\char`\}, \char`\{$x_2$,$y_2$\char`\}, ...\char`\}, \char`\{$f_1(x)$, 
$f_2(x)$, ...\char`\}, $x$]}
\item{} where the first argument contains a list of pair of points,
the second argument is a list of linearly independent basis 
functions, and the last argument (a symbol) is the independent variable of the functions. 
A polynomial fit is obtained if we take $f_i(x) = x^i$, $i = 0, 1, \ldots, n$.
The function returns the fitted result
$$F(x) = \sum_i c_i f_i(x)$$
\item{} such that the square deviation of error is minimized.  For some more
discussions, see T.~H.~Cormen, C.~E.~Leiserson, R.~L.~Rivest, ``Introduction
to Algorithms,'' p. 768--771.  We use the result
that the fitting coefficient $c_i$ taken as a column vector
is determined by the equation
$$ A^{T}A c = A^{T} y, $$
\item{} where the matrix $A$ is defined by $A_{ij} = f_j(x_i)$, 
(the dimensions of $A$ are the number of sample points by the number of 
fitting functions) and $A^{T}$ is the
transpose of $A$, and $y$ is the column vector $y_i$.   Use the matrix
and linear algebra capability of Mathematica in your definition.
Given some testing sample outputs (e.g. a plot with points and fitted curve).

\problem  [A difficult one] (a) Write a function {\tt buildtree[]} that
takes a list and builds a binary tree with nested lists.  
Each node is represented as  
\ttbg
{ nodeelement, {leftchild}, {rightchild} }
\tted
\item{}The children in turn have similar recursive data structure 
representation.  The order of the elements in the tree is such that 
all the nodes of left subtree is not late in the
Mathematica internal ordering than the parent, as determined 
by {\tt Sort[ ]}. 
The elements of the input list are used from the beginning one
by one, and are added to the tree.  Here are some examples,
\ttbg
buildtree[{5,3,2,9}]  => {5, {3, {2, {}, {}}, {}}, {9, {}, {}}}
buildtree[{5,c,2,a}]  => {5, {2, {}, {}}, {c, {a, {}, {}}, {}}}
buildtree[{{a,b},{1,2},{u,v,w}}]=>{{a,b},{{1,2},{},{}},{{u,v,w},{},{}}}
\tted
\item{} Note that the nodes themselves can be arbitrary expressions. 
{\sl Hint}: Define a separate function {\tt addelement} for adding
one element to a tree.  Think in terms of recursive call, but do not think in
terms of pointers.  There is no pointer in Mathematica. 
\item{} (b) Write a function {\tt traverse[]} that takes a tree
and builds a list in-order; that is, in an order
that visit the left subtree first, then the node, and then the right
subtree.  {\tt traverse[buildtree[$e$]]} produces the same
result as {\tt Sort[$e$]}.  here are three examples:
\ttbg
traverse[buildtree[{5,3,2,9}]]  => {2, 3, 5, 9}
traverse[buildtree[{5,c,2,a}]]  => {2, 5, a, c}
traverse[buildtree[{{a,b},{1,2},{u,v,w}}]] => {{1,2},{a, b},{u,v,w}}
\tted
Some background in computer algorithms will help.

\bye
